\documentclass[10pt]{article}
\usepackage[margin=0.78in]{geometry}
\usepackage{booktabs}
\usepackage{graphicx}
\usepackage{hyperref}
\usepackage{xcolor}
\usepackage{float}
\usepackage{caption}
\hypersetup{colorlinks=true,linkcolor=blue!55!black,citecolor=blue!55!black,urlcolor=blue!55!black}
\title{First-in Decisions at Ten Big Blinds in PLO Tournaments:\\
A Restricted-Action Pilot Across Four Hand Archetypes}
\author{}
\date{September 2026}

\begin{document}
\maketitle

\begin{abstract}
\input{generated/abstract_results.tex}
\end{abstract}

\section{Question and hypotheses}
This pilot asks how position and tournament stage change a first-in decision with a ten-big-blind
stack in six-handed pot-limit Omaha high. The decision is deliberately restricted to fold, limp to
one big blind, or raise to two big blinds. A pot-sized opening raise is outside the model, while
pot-sized three-bets, later preflop raises, and postflop pot bets remain available. The analysis is
therefore about one-step action values inside a specified abstraction. It is not a claim about a
complete PLO strategy or a Nash equilibrium.

The primary prespecified user hypothesis is that Premium and Speculative examples may favor
limping, while Marginal examples may favor folding; the Trash examples provide an additional
weak-hand comparison. Position and the 76-versus-60-player snapshots are exploratory contrasts,
without a directional prediction. In particular, SB is later in the preflop order but remains out
of position after the flop, so a monotone seat effect is not assumed. These hypotheses are not
encoded action rules. Eight selected hands cannot identify population-wide frequencies for an
archetype.

\section{Tournament and game assumptions}
The tournament profile reproduces the Sunday Classic Mini example: 515 initial entries, 75 paid,
and the user-supplied ordered prize schedule encoded in the configuration, totaling \$12,500. We
study 76 and 60 players remaining as separate hypothetical snapshots, not sequential states of one
simulated tournament. At 60 remaining, the encoded top-60 prizes sum to \$11,684.90; the \$815.10
assigned to ranks 61--75 is no longer available. Absolute EV levels therefore should not be compared
across snapshots as if the payout pools were identical; action contrasts within each snapshot are
the primary quantities.
All six players at the table and all players outside it start each modeled hand with 10 big blinds
before forced posts. Every seated player posts a 0.1-big-blind ante, the small blind is 0.5, the big
blind is 1, and rake is zero. The six seats are UTG, HJ, CO, BTN, SB, and BB; the focal first-in
decisions cover UTG through SB.

The first-in ``limp to one big blind'' is a total preflop street contribution: UTG through BTN add
1 big blind, while SB adds 0.5 after its blind. A raise to 2 big blinds similarly makes SB add 1.5.
Antes are separate sunk contributions, and fold values retain the effect of forced posts. BB is not
an RFI study seat because the hand ends when every other live player folds to it.

The engine plays all four streets. It enforces pot-limit sizing, short all-ins and reopening,
side pots, ties, and Omaha's exactly-two-hole-card and exactly-three-board-card showdown rule.
Antes are excluded from the preflop GGPoker pot-limit cap and included on later streets. In an
unraised preflop pot the first raise abstraction is exactly two big blinds, clipped to the player's
actual all-in amount. After a raise, the two-big-blind opening action disappears and pot-sized
reraises remain. Postflop actions are check, call, fold, or pot where legal.

\pagebreak[3]
\section{Hand cohort}
The cohort contains two exact physical hands in each of four locally defined drill tiers:
\begin{center}
\begin{tabular}{lll}
\toprule
Tier & Hand 1 & Hand 2 \\
\midrule
Premium & A\(\spadesuit\)A\(\heartsuit\)K\(\spadesuit\)K\(\heartsuit\) & J\(\spadesuit\)T\(\clubsuit\)9\(\spadesuit\)8\(\clubsuit\) \\
Speculative & A\(\spadesuit\)A\(\clubsuit\)8\(\diamondsuit\)2\(\heartsuit\) & A\(\diamondsuit\)9\(\diamondsuit\)8\(\spadesuit\)6\(\spadesuit\) \\
Marginal & J\(\diamondsuit\)J\(\clubsuit\)6\(\clubsuit\)3\(\spadesuit\) & K\(\diamondsuit\)J\(\heartsuit\)T\(\diamondsuit\)9\(\spadesuit\) \\
Trash & 9\(\heartsuit\)7\(\diamondsuit\)5\(\clubsuit\)3\(\spadesuit\) & Q\(\spadesuit\)J\(\spadesuit\)7\(\clubsuit\)6\(\clubsuit\) \\
\bottomrule
\end{tabular}
\end{center}
The software reclassifies each exact hand before a run and aborts if its tier differs from the
frozen cohort. AAKK is the user's classifier-derived anchor; the other seven labels are cached drill
examples in \path{harnesses/EN/sources/web/plo-starting-hands-classify-*}. These tier names use
Hwang-style starting-hand structure. They select varied examples and are not measured latent hand
strength. Tier labels are used only to validate and report the cohort; they are never supplied to
the solver as action rules or rewards.

\section{Computation}
\subsection{Abstract policy training}
The trainer uses external-sampling Monte Carlo counterfactual regret minimization (MCCFR), following
the sampling framework of Lanctot et al.~\cite{lanctot2009}. Six seats train in the full multiplayer
tree. Private observations use finite structural buckets for rank multiplicity, aces, suit pattern,
high-card count, coarse connectivity, pair band, the PokerKit made-hand category, board rank and suit
multiplicities, and combined hole-plus-board suit counts. They do not explicitly encode wrap outs,
nut-flush status, blockers, or exact kickers. Keys
retain the complete public action history, public stacks, and the focal player's own sequence of
observation buckets. Opponent keys never include the focal query flag or its private cards.

Chance sampling mixes ordinary deals with the global-suit orbit of the queried cohort representative.
Thus a run preserves the rank and suit-isomorphism class rather than always dealing the same named
physical suits. Regret
and average-strategy updates receive the associated importance weight once. Strategy averaging uses
a separate full-support sampling pass. Each main cell uses a 15,000-node stopping threshold (actual
counters are archived), one of three seeds (17, 29, 43), and a frozen source snapshot. No
convergence claim follows from this finite
budget or from MCCFR theory for the present multiplayer, abstracted game.

\subsection{One-step action values}
After training, policies are frozen. Evaluation deals the focal hand's suit-isomorphism class independently and samples
the other private cards and future board without replacement. For positions after UTG, a deal is
weighted by the product of learned average fold probabilities for preceding seats. Each legal focal
action is then forced on a common deal and common random continuation. All remaining players and
streets continue under the frozen average policies; folding the focal player does not terminate the
hand prematurely. If a preceding or continuation information bucket was unvisited, its explicit
uniform fallback policy contributes to both the fold weight and the rollout; the fallback fraction
is reported. Specifically, it is the unweighted fraction of policy lookups using uniform fallback
during prefix conditioning and completed paired continuation evaluation, including prefix lookups
on attempts whose conditioning weight becomes zero. It is neither posterior-weighted nor a played
strategy frequency.

For action \(a\), the reported value is the weighted mean
\[
  \widehat V_a=\frac{\sum_j w_j U_{ja}}{\sum_j w_j}, \qquad
  \mathrm{ESS}=\frac{(\sum_j w_j)^2}{\sum_j w_j^2}.
\]
Paired standard errors are computed from within-deal differences, preserving covariance between
actions. We report raise-minus-limp and limp-minus-fold contrasts directly. The chosen action is the
largest estimated mean in that cell. It is one-step selection under frozen continuation policies,
not a full best response. Each main cell requests 256 evaluation attempts. An interrupted final
pair is discarded and marks the cell truncated.

\subsection{ICM conversion and reproducibility}
Terminal chip stacks are converted to payout dollars through the independent chip model (ICM).
The implementation uses the exact Harville recurrence and compresses identical outside-field
stacks into a counted crowd. For simultaneous busts it orders tied players by their pre-hand
stacks; equal pre-hand stacks split the corresponding prizes. Thus terminal ICM valuation adds no
Monte Carlo noise.
Gilbert formalizes risk effects within ICM~\cite{gilbert2009}; recent empirical work also documents
systematic limits in how ICM represents tournament outcomes~\cite{kim2025}. We therefore treat ICM
as the declared valuation model rather than ground truth. Poker hand settlement uses PokerKit's
Omaha evaluator~\cite{kim2023pokerkit}.

The main grid contains 8 hands \(\times\) 5 positions \(\times\) 2 stages \(\times\) 3 seeds,
or 240 cells. Every cell stores its configuration, source hash, exact cards, training counters,
three action values, all paired standard errors, effective sample size, fallback fraction, runtime,
and completion status. Results from different seeds are summarized as a mean and range, while modal
action agreement is shown as a count out of three. These are separate trained policies and are not
pooled into a single confidence interval. Every cell retrains its own query-specific model; the
position figures do not represent one shared, coherent opening policy.

\section{Results}
\input{generated/results_text.tex}

\begin{figure}[H]
\centering
\includegraphics[width=\textwidth]{generated/modal_actions_left76.pdf}
\caption{Modal one-step action by exact hand and position with 76 players remaining. Parentheses
show modal votes among completed cells and $n$ reports completed seeds out of three. Gray cells
indicate missing results.}
\end{figure}

\begin{figure}[H]
\centering
\includegraphics[width=\textwidth]{generated/modal_actions_left60.pdf}
\caption{The corresponding positional summary with 60 players remaining. Colors encode the modal
action only; they are not mixed-strategy frequencies.}
\end{figure}

\input{generated/position_tables.tex}

\section{Budget sensitivity}
\input{generated/sensitivity_text.tex}

\section{Limitations}
\begin{enumerate}
\item The opening action set excludes pot-sized opens by design. The study cannot show that a
pot-sized open is inferior; it only compares fold, limp, and raise-to-two within this model. The
continuation abstraction also omits smaller postflop bet and raise sizes, which can change preflop
action values and rankings.
\item The private and postflop information abstraction pools strategically different hands. Uniform
fallback at unvisited decision points is explicit and is reported as a diagnostic, not learned play.
\item Eight exact hands are selected examples. Two hands per named tier do not support estimates of
the distribution of all PLO hands in that tier.
\item Every hand--position--stage--seed cell retrains a focal-hand abstract policy. Within a cell,
all three focal actions share one frozen continuation policy and common deals. Across cells, the
trained opponent policies can differ. The grid is therefore a collection of local action-value
experiments and cannot be assembled into one coherent equilibrium opening range.
\item This is a six-player general-sum tournament setting with finite MCCFR work. The reported
policies have no exploitability bound and no Nash or GTO guarantee.
\item Sampling standard errors quantify rollout variation conditional on a frozen trained policy.
They exclude abstraction error, ICM model error, training uncertainty, and selection bias from taking
the largest estimated action value.
\item Equal 10-big-blind stacks outside the table and the exact prize schedule materially determine
ICM values. Different field stacks, payouts, rake, antes, or prior action define a different problem.
The single-hand terminal valuation does not simulate busts at other tables, future blind increases,
future hands, or skill differences. Equal outside stacks also remove micro-stack ``wait out the
bubble'' pressure that can materially change a real decision.
\end{enumerate}

\pagebreak[3]
\section{Conclusion}
\input{generated/conclusion_text.tex}

\appendix
\section{Encoded payout schedule}
\begin{center}
\begin{tabular}{rrr}
\toprule
Finishing rank(s) & Places & Prize per place \\
\midrule
1 & 1 & \$2,087.06 \\
2 & 1 & \$1,547.88 \\
3 & 1 & \$1,148.27 \\
4 & 1 & \$851.83 \\
5 & 1 & \$631.92 \\
6 & 1 & \$468.78 \\
7 & 1 & \$347.75 \\
8 & 1 & \$236.87 \\
9--10 & 2 & \$191.94 \\
11--12 & 2 & \$155.54 \\
13--16 & 4 & \$126.04 \\
17--22 & 6 & \$102.13 \\
23--32 & 10 & \$82.76 \\
33--48 & 16 & \$67.06 \\
49--75 & 27 & \$54.34 \\
\bottomrule
\end{tabular}
\end{center}

\section*{Reproduction}
The frozen run records its source snapshot and manifest beside cell JSON. From the repository root:
{\scriptsize
\begin{verbatim}
.venv/bin/python research/plo_icm_archetypes/run_study.py run \
  --output tmp/plo_icm/study_main_v1 --max-nodes 15000 \
  --eval-attempts 256 --timeout 60 --wall-limit 2700 --workers 2
.venv/bin/python tmp/plo_icm/study_main_v1/source_snapshot/research/plo_icm_archetypes/run_study.py \
  sensitivity --output tmp/plo_icm/study_sensitivity_v1 \
  --max-nodes 60000 --eval-attempts 1024 --timeout 180 \
  --wall-limit 720 --workers 2
.venv/bin/python research/plo_icm_archetypes/analyze_results.py \
  --main tmp/plo_icm/study_main_v1 \
  --sensitivity tmp/plo_icm/study_sensitivity_v1 \
  --output research/plo_icm_archetypes/generated
\end{verbatim}
}

\clearpage
\begin{thebibliography}{9}
\bibitem{lanctot2009}
M. Lanctot, K. Waugh, M. Zinkevich, and M. Bowling.
Monte Carlo sampling for regret minimization in extensive games.
\emph{Advances in Neural Information Processing Systems 22}, 2009.
\url{https://papers.nips.cc/paper/3713-monte-carlo-sampling-for-regret-minimization-in-extensive-games}

\bibitem{gilbert2009}
G. T. Gilbert. The Independent Chip Model and Risk Aversion. arXiv:0911.3100, 2009.
\url{https://arxiv.org/abs/0911.3100}

\bibitem{kim2025}
J. Kim. Empirical Validation of the Independent Chip Model. arXiv:2506.00180, 2025.
\url{https://arxiv.org/abs/2506.00180}

\bibitem{kim2023pokerkit}
J. Kim. PokerKit: A Comprehensive Python Library for Fine-Grained Multi-Variant Poker Game
Simulations. arXiv:2308.07327, 2023.
\url{https://arxiv.org/abs/2308.07327}
\end{thebibliography}
\end{document}
